\chapter{Introducción}\label{cap:introduccion}

\section{Planteamiento del problema}

La detección de lavado de dinero se apoya cada vez más en modelos de aprendizaje automático
que examinan transacciones en busca de patrones anómalos \cite{Chen2018AMLReview,
Kute2021AMLReview}. En el caso de las criptomonedas, las transacciones forman de manera natural
un grafo, en el que cada transacción es un nodo y cada transferencia de fondos entre
transacciones es una arista. Las redes neuronales de grafos (\textit{Graph Neural Networks},
GNN) aprovechan esa estructura combinando los atributos de cada transacción con los de sus
vecinas \cite{Scarselli2009GNN, Wu2021Survey}, y se han convertido en una de las familias de
modelos más estudiadas para la detección de fraude financiero \cite{Motie2024FraudGNN}.

En un entorno regulado no basta con que el modelo acierte. El analista que revisa una alerta
necesita saber por qué la transacción fue marcada como sospechosa, y esa justificación debe
poder sostenerse ante un auditor \cite{Rudin2019Stop, Arrieta2020XAI}. Las GNN son modelos
opacos, por lo que la justificación se obtiene con métodos de explicabilidad
post hoc que asignan una importancia a cada variable o a cada arista que intervino en la
predicción \cite{Yuan2023Taxonomic}. Esos métodos, sin embargo, son en su mayoría
estocásticos u optimizan una máscara a partir de una inicialización aleatoria. Si se ejecutan
dos veces sobre la misma transacción y el mismo modelo, pueden devolver explicaciones
distintas. Una explicación que cambia cada vez que se pide no sirve como evidencia.

A esa fuente de variación se suma otra, propia del dominio. Las transacciones ilícitas son
una fracción muy pequeña del total. En el conjunto de datos Elliptic, que es la referencia
pública más usada para este problema \cite{Weber2019Elliptic}, hay 4.545 transacciones
etiquetadas como ilícitas frente a más de 199.000 lícitas o sin etiqueta. Bajo un desbalance
así, los modelos aprenden la clase rara a partir de pocos ejemplos \cite{He2009Imbalanced,
Johnson2019Survey}, y es razonable esperar que sus explicaciones sobre esa clase también sean
más frágiles. Las técnicas que compensan el desbalance, como reponderar la función de
pérdida o generar ejemplos sintéticos de la clase rara \cite{Chawla2002SMOTE, Lin2020Focal},
modifican lo que el modelo aprende y podrían modificar también lo que el explicador reporta.

La literatura ha evaluado la calidad de los explicadores de GNN principalmente en grafos
sintéticos o en tareas de química y redes sociales \cite{Agarwal2023Evaluating,
Longa2024Explaining}, y la evaluación de la estabilidad de las explicaciones se ha estudiado
sobre todo fuera de los grafos \cite{Nauta2023Anecdotal, Bodria2023Benchmarking}. No se
encontró un estudio que mida, con un diseño factorial y réplicas independientes, cómo cambian
la estabilidad de las explicaciones de distintas GNN cuando se varían el nivel de desbalance y
la técnica que lo compensa, en un problema real de lavado de dinero. Ese es el vacío que
aborda esta tesis.

\section{Una limitación del conjunto de datos que condiciona el diseño}

Elliptic cubre 49 pasos de tiempo separados aproximadamente dos semanas entre sí. En el paso
43 se cerró un mercado ilegal importante de la red, y los autores del conjunto de datos
reportaron que, a partir de ese momento, ningún modelo era capaz de detectar las nuevas
transacciones ilícitas \cite{Weber2019Elliptic}. La partición temporal habitual entrena con
los pasos 1 a 34, valida con los pasos 35 a 42 y evalúa con los pasos 43 a 49, de modo que la
prueba cae completa después del cierre. En la primera versión de los experimentos de esta
tesis, todos los modelos quedaron en el nivel del azar en esa prueba (Anexo
\ref{anx:original}).

Por esa razón, el estudio principal usa solo los pasos 1 a 42, anteriores al cierre, y los
reparte en entrenamiento (1 a 29), validación (30 a 35) y prueba (36 a 42). Los pasos 43 a 49
no se usan. La decisión, sus costos y la comparación con la partición habitual se exponen en
el Capítulo \ref{cap:metodologia}.

\section{Preguntas de investigación}

\textbf{Pregunta principal.} ¿Cómo cambia la estabilidad de las explicaciones que producen
los métodos de explicabilidad sobre redes neuronales de grafos entrenadas para detectar
transacciones ilícitas en Elliptic, cuando se varía el nivel de desbalance de clases, la
arquitectura de la red y la técnica empleada para compensar el desbalance?

\textbf{Preguntas específicas.}
\begin{enumerate}
  \item ¿Disminuye la estabilidad de las explicaciones a medida que el desbalance de clases
        con el que se entrena el modelo se hace más extremo?
  \item ¿Producen las distintas arquitecturas (GCN, GraphSAGE, GAT y TAGCN) explicaciones con
        distinta estabilidad?
  \item ¿Cambia la estabilidad de las explicaciones según la técnica usada para compensar el
        desbalance (sin ajuste, pesos por clase o \textit{focal loss})?
  \item ¿Refleja la estabilidad medida lo que el modelo aprendió, o se debe principalmente a
        los datos de entrada y al propio explicador?
\end{enumerate}

\section{Hipótesis}\label{sec:hipotesis}

Las hipótesis se formularon antes de observar los resultados y se contrastan en el Capítulo
\ref{cap:resultados}.

\begin{description}
  \item[H1 (desbalance).] La estabilidad de las explicaciones disminuye a medida que el
        desbalance de clases del entrenamiento se hace más extremo, porque el modelo dispone de
        menos ejemplos de la clase ilícita de los cuales aprender.
  \item[H2 (arquitectura).] La estabilidad de las explicaciones depende de la arquitectura,
        porque cada una combina de forma distinta la información de la transacción con la de
        sus vecinas.
  \item[H3 (compensación).] La técnica de compensación del desbalance no cambia la
        estabilidad de las explicaciones, porque modifica cuánto pesa cada clase en la
        función de pérdida pero no la información disponible para el modelo.
\end{description}

\section{Objetivos}

\subsection{Objetivo general}

Evaluar la estabilidad de los métodos de explicabilidad aplicados a redes neuronales de grafos
entrenadas para detectar lavado de dinero en el conjunto de datos Elliptic, bajo distintos
niveles de desbalance de clases, arquitecturas y técnicas de compensación, con el fin de
establecer qué factores la determinan y en qué medida sus explicaciones son confiables para
una auditoría.

\subsection{Objetivos específicos}

\begin{enumerate}
  \item Medir el efecto del nivel de desbalance de clases, desde 1:10 hasta 1:200, sobre la
        estabilidad de las explicaciones de cinco explicadores: GNNExplainer, PGExplainer,
        valores de Shapley por muestreo, Integrated Gradients y Expected Gradients.
  \item Comparar la estabilidad de las explicaciones entre cuatro arquitecturas de GNN: GCN,
        GraphSAGE, GAT y TAGCN.
  \item Analizar el efecto de tres técnicas de compensación del desbalance (sin ajuste, pesos
        por clase y \textit{focal loss}) y del sobremuestreo con SMOTE sobre el rendimiento
        predictivo y sobre la estabilidad de las explicaciones.
  \item Determinar, mediante un control con pesos aleatorios y medidas de acuerdo que no
        dependen de las variables nulas, qué parte de la estabilidad medida proviene del modelo
        entrenado y qué parte de los datos de entrada y del explicador.
\end{enumerate}

Respecto de la propuesta aprobada, dos métodos se sustituyeron por alternativas que se ajustan
mejor al problema. GNNShap se reemplazó por valores de Shapley
estimados por muestreo sobre las variables de la transacción \cite{Strumbelj2014Explaining},
porque la pregunta de la tesis es qué variables pesan y GNNShap atribuye importancia a las
aristas. GraphSMOTE se reemplazó por SMOTE sobre los atributos \cite{Chawla2002SMOTE}, porque
la estructura del grafo de Elliptic no se modifica en ningún escenario. Integrated Gradients
\cite{Sundararajan2017IG} y Expected Gradients \cite{Erion2021Expected} se añadieron como
explicadores de referencia, uno determinista y uno estocástico, que no requieren entrenamiento.

\section{Contribuciones}

\begin{enumerate}
  \item Un diseño factorial replicado (4 arquitecturas, 6 escenarios de desbalance, 3 técnicas
        de compensación y 3 semillas de modelo, 216 modelos) que mide la estabilidad de cinco
        explicadores sobre las mismas 30 transacciones ilícitas.
  \item La evidencia de que la partición temporal habitual de Elliptic deja la prueba
        completa después del cierre del mercado ilegal y vuelve inútil la evaluación en
        prueba, junto con una partición alternativa que la recupera.
  \item La identificación de un artefacto de medida: cerca de la mitad de las variables de
        una transacción valen cero después del escalado, y los explicadores las dejan siempre
        al final del orden de importancia, lo que infla la correlación de Spearman entre
        explicaciones. La tesis cuantifica ese efecto con dos medidas que no dependen de esas
        variables y con un control de pesos aleatorios \cite{Adebayo2018Sanity}.
\end{enumerate}

\section{Alcance y limitaciones}

El estudio se limita al conjunto de datos Elliptic y a la tarea de clasificar transacciones
como ilícitas o no ilícitas. Las transacciones sin etiqueta se tratan como no ilícitas, que es
el supuesto más cercano a la operación real (Capítulo \ref{cap:metodologia}). La estabilidad se
mide sobre explicaciones de variables, salvo en PGExplainer, que solo explica aristas. No se
evalúa la plausibilidad de las explicaciones frente a un patrón real de lavado, porque Elliptic
no publica el significado de sus variables ni anotaciones de tipologías; un grafo sintético
construido con ese fin se resume en el Anexo \ref{anx:sintetico}. No se contempla el despliegue
del sistema ni su evaluación con analistas.

\section{Estructura del documento}

El Capítulo \ref{cap:marco} presenta los fundamentos de las GNN, del desbalance de clases y de
los métodos de explicabilidad usados. El Capítulo \ref{cap:estado} revisa los trabajos sobre
detección de lavado con grafos y sobre la evaluación de explicaciones. El Capítulo
\ref{cap:metodologia} describe los datos, la partición temporal, los escenarios, el
entrenamiento, los explicadores, las medidas de estabilidad y las pruebas estadísticas. El
Capítulo \ref{cap:resultados} presenta los resultados y contrasta las hipótesis. El Capítulo
\ref{cap:discusion} interpreta los resultados y sus límites, y el Capítulo
\ref{cap:conclusiones} resume las respuestas a las preguntas de investigación. Los anexos
recogen la corrida con la partición original, el eje sintético y los detalles de
reproducibilidad.
